%% kapa: exam=1 round=35 year=2024 period=2 form=A subject=accounting subject_ko=회계학 seq=2 number=42
%% source: Q-Net 감정평가사 기출문제 (cst00309) artlSeq=5213229 file=raw/1008/5213229/2교시 문제지 원본.pdf
%% conversion: kordoc 4.24.51 -> tools/gen1.py (auto); review=proofread-sonnet; flags=none
\begin{question}{42}
재무제표 표시에 관한 설명으로 옳은 것은?
\begin{choices}
  \item 기업이 재무상태표에 유동자산과 비유동자산, 그리고 유동부채와 비유동부채로 구분하여 표시하는 경우, 이연법인세자산은 유동자산으로 분류한다.
  \item 한국채택국제회계기준을 준수하여 작성된 재무제표는 국제회계기준을 준수하여 작성된 재무제표임을 주석으로 공시할 수 있다.
  \item 환경 요인이 유의적인 산업에 속해 있는 경우나 종업원이 재무제표이용자인 경우 재무제표 이외에 환경보고서나 부가가치보고서도 한국채택국제회계기준을 적용하여 작성한다.
  \item 부적절한 회계정책은 이에 대하여 공시나 주석 또는 보충자료를 통해 설명하여 정당화될 수 있다.
  \item 당기손익과 기타포괄손익은 별개의 손익계산서가 아닌 단일의 포괄손익계산서로 작성되어야 한다.
\end{choices}
\end{question}
\answer{2}
